% Ampliación sustantiva de los comparadores determinantes BSD.
% Módulo destinado al capítulo BSD de la Parte IV.

\section{Comparateurs de l'unité déterminantale commune}
\label{sec:amp-bsd-comparadores}

L'égalité racine \eqref{eq:bsd-suc-raiz} fixe une seule unité dans les canaux
Kato et Poitou--Tate. À partir de cette unité, on construit maintenant les
comparateurs qui transportent le régulateur dérivé jusqu'à sa publication
scalaire. Aucun comparateur ne choisit une normalisation après avoir connu le
terme principal : tous proviennent de morphismes de complexes, de suites
exactes et de bases déjà orientées.

\subsection{La chaîne de droites déterminantes}
\label{subsec:amp-bsd-lineas}

Pour un complexe parfait basé \(C\), on utilise la convention
\begin{equation}\label{eq:amp-bsd-linea-determinante}
 \lambda(C):=\bigotimes_i
 \det H^i(C)^{(-1)^i}.
\end{equation}
Soient \(\mathcal C_E^K\), \(\mathcal C_E^{\mathrm{Iw}}\),
\(\mathcal C_E^{\mathrm{PT}}\) et \(\mathcal C_E^{\mathrm{loc}}\),
respectivement, le complexe publié par le canal de Kato, son modèle
d'Iwasawa basé, le complexe réduit autodual de Poitou--Tate et le cône de
localisation finie--singulière. Leurs droites sont reliées par
\begin{equation}\label{eq:amp-bsd-cadena-lineas}
 \mathcal L_E^K
 \xrightarrow{\ \mathfrak c_{K/\mathrm{FERS}}\ }
 \mathcal L_E^{\mathrm{Iw}}
 \xrightarrow{\ \mathfrak c_{\mathrm{PT}}\ }
 \mathcal L_E^{\mathrm{ht}}
 \xrightarrow{\ \mathfrak c_{\mathrm{STS}}\ }
 \mathcal L_E^{\mathrm{fin}}
 \xrightarrow{\ \mathfrak c_{\mathrm{tors}}\ }
 \mathcal L_E^{\mathrm{ar}}
 \xrightarrow{\ \mathfrak c_{\infty}\ }
 \mathcal L_E.
\end{equation}
Le sigle \(\mathrm{STS}\) désigne dans cette formule l'étape conjointe de
Smith, Tamagawa et Tate--Shafarevich. Toutes les flèches de
\eqref{eq:amp-bsd-cadena-lineas} sont des isomorphismes de droites orientées.

\subsection{Comparaison de chaînes Kato--FERS à partir de l'unité commune}
\label{subsec:amp-bsd-comparacion-cadenas-kato-fers}

L'équivalence à l'origine du premier comparateur peut être construite avant de
prendre les déterminants. Soit \(F^q\mathcal C\) la filtration par profondeur
nonadique, identifiée à la filtration \(T\)-adique au moyen de la retenue, et
écrivons
\begin{equation}\label{eq:amp-bsd-complejos-kato-iwasawa}
 \mathcal C_E^K
 :=\operatorname{Tot}\!\left(
 P_E^K\widehat\Pi_0G_{m,n}\right),
 \qquad
 \mathcal C_E^{\mathrm{Iw}}
 :=[V_\Lambda\xrightarrow{S_E(T)}W_\Lambda].
\end{equation}
Soit
\(\operatorname{car}_T:G_{m,n}\to\mathcal C_E^{\mathrm{Iw}}\) le lecteur
de retenue qui associe à une route \(\gamma\) le générateur déterminé par ses
faces initiale et finale, multiplié par
\(T^{\operatorname{Carr}(\gamma)}\). La multiplication des concaténations dans
le modèle d'Iwasawa est notée \(\mu_{\mathrm{Iw}}\). La diagonale
d'Alexander--Whitney et la publication \(P_E^K\) définissent alors, avant de
prendre la cohomologie ou les déterminants, l'application de chaînes
\begin{equation}\label{eq:amp-bsd-phi-kato-fers}
\begin{aligned}
 \Phi_{K/\mathrm{FERS}}&:
 \mathcal C_E^K\longrightarrow\mathcal C_E^{\mathrm{Iw}},\\
 \Phi_{K/\mathrm{FERS}}
 &:=
 \mu_{\mathrm{Iw}}\circ
 \bigl(P_E^K\widehat\otimes\operatorname{car}_T\bigr)
 \circ\Delta_{\mathrm{AW}},\\
 \Phi_{K/\mathrm{FERS}}(z_K)&=z_{\mathrm{Iw}}.
\end{aligned}
\end{equation}
où \(z_{\mathrm{Iw}}\) est le générateur racine induit par \(1_E\). Ainsi,
chaque simplexe normalisé est envoyé sur la somme de ses coupes
d'Alexander--Whitney, en conservant dans chaque terme les deux faces et la puissance
de retenue. La formule ne consulte pas le terme principal analytique.

\begin{lemma}[Critère générateur par générateur pour l'équivalence]
\label{lem:amp-bsd-bases-emparejadas-kato-fers}
À chaque degré \(i\), soit \(\mathcal B_i^K\) la base finie de simplexes
normalisés, ordonnée par profondeur, faces et retenue, et soit
\(\mathcal B_i^{\mathrm{Iw}}\) la base formée par les générateurs ayant les
mêmes faces initiale et finale dans le modèle d'Iwasawa. Pour une application
candidate
\[
 \upsilon_i:\mathcal B_i^K\longrightarrow
 \mathcal B_i^{\mathrm{Iw}},
\]
définissons le résiduel de degré zéro
\begin{equation}\label{eq:amp-bsd-residual-generador-graduado}
 \varepsilon_i(b)
 :=\Phi_{K/\mathrm{FERS}}^i(b)\bmod T-\upsilon_i(b).
\end{equation}
Si \(\upsilon_i\) est une bijection et
\(\varepsilon_i(b)=0\) pour tout \(b\in\mathcal B_i^K\), alors les deux
modules complétés de degré \(i\) sont libres de même rang
\(d_i=|\mathcal B_i^K|\) et, pour chaque générateur,
\begin{equation}\label{eq:amp-bsd-phi-en-cada-generador}
 \Phi_{K/\mathrm{FERS}}^i(b)
 =\upsilon_i(b)
  +T\sum_{c\in\mathcal B_i^{\mathrm{Iw}}}
    \lambda_{c,b}(T)c,
 \qquad \lambda_{c,b}(T)\in\Lambda.
\end{equation}
En particulier,
\(\operatorname{gr}^{F}\Phi_{K/\mathrm{FERS}}^i=I\) sur tous les
générateurs, et non seulement sur l'unité racine.
\end{lemma}

\begin{proof}
L'annulation de \eqref{eq:amp-bsd-residual-generador-graduado} signifie que la
colonne de \(b\) a pour terme constant \(\upsilon_i(b)\). Les autres
coefficients appartiennent à \(T\Lambda\), ce qui donne
\eqref{eq:amp-bsd-phi-en-cada-generador}. La bijectivité identifie les
bases et fixe le même rang. Dans ces bases, la réduction modulo \(T\) de la
matrice de \(\Phi^i\) est l'identité ; d'où
\(\operatorname{gr}^{F}\Phi^i=I\).
\end{proof}

L'égalité racine
\(\Phi_{K/\mathrm{FERS}}(z_K)=z_{\mathrm{Iw}}\) fixe la normalisation commune.
L'écriture de \(P_E^K\) sur chaque simplexe normalisé, la bijection
\(\upsilon_i\) et les résiduels \(\varepsilon_i(b)\) fournissent un
contrôle générateur par générateur de l'équivalence déjà construite ; ils ne choisissent ni ne
régissent la conclusion BSD.

\begin{lemma}[Équivalence filtrée et normalisation de la racine]
\label{lem:amp-bsd-equivalencia-kato-fers}
Pour l'application de chaînes construite dans
\eqref{eq:amp-bsd-phi-kato-fers}, l'application \(\upsilon_i\) est la
correspondance de bases induite par les faces et la retenue de chaque
simplexe. La formule d'Alexander--Whitney vérifie à chaque degré le
critère du lemme \ref{lem:amp-bsd-bases-emparejadas-kato-fers} ; il s'agit d'un
contrôle de l'équivalence déjà construite, et non d'une prémisse supplémentaire.
L'application \(\Phi_{K/\mathrm{FERS}}\) est une équivalence basée de
complexes parfaits. Dans les bases ordonnées par profondeur, ses matrices à
chaque degré sont de la forme
\begin{equation}\label{eq:amp-bsd-matriz-kato-fers}
 [\Phi_{K/\mathrm{FERS}}^i](T)=I+TB_i(T),
 \qquad B_i(T)\in M_{d_i}(\Lambda).
\end{equation}
Par conséquent, l'unité
\begin{equation}\label{eq:amp-bsd-rho-determinantal}
 \rho_{E,\ell}(T)
 :=\prod_i\det\!\left(I+TB_i(T)\right)^{(-1)^i}
 \in\Lambda^\times
\end{equation}
satisface \(\rho_{E,\ell}(0)=1\).
\end{lemma}

\begin{proof}
L'identité simpliciale
\(\partial\Delta_{\mathrm{AW}}
=(\partial\otimes I+(-1)^{\deg}I\otimes\partial)
\Delta_{\mathrm{AW}}\), avec la compatibilité de
\(P_E^K\) avec les faces et les dégénérescences, démontre que
\eqref{eq:amp-bsd-phi-kato-fers} commute avec les différentielles. Dans le gradué
associé, le lemme
\ref{lem:amp-bsd-bases-emparejadas-kato-fers} identifie tous les
générateurs et donne l'identité. Dans ces bases,
\eqref{eq:amp-bsd-phi-en-cada-generador} est précisément
\eqref{eq:amp-bsd-matriz-kato-fers}.

La complétude de \(\Lambda=K[[T]]\) fait converger la série
\begin{equation}\label{eq:amp-bsd-inversa-kato-fers}
 (I+TB_i(T))^{-1}
 =\sum_{n\geq0}(-TB_i(T))^n;
\end{equation}
ces inverses commutent également avec les différentielles, et forment donc
l'inverse basé de \(\Phi_{K/\mathrm{FERS}}\). En prenant le déterminant
alterné, on obtient \eqref{eq:amp-bsd-rho-determinantal} ; l'évaluation en
\(T=0\) donne un. Ainsi, la normalisation n'est pas postulée à partir du terme
principal : c'est la coordonnée constante d'une équivalence qui est l'identité
sur l'unité racine.
\end{proof}

Le premier comparateur est le déterminant de
\(\Phi_{K/\mathrm{FERS}}\). Si \(\mathfrak z_K\) désigne l'élément
provenant de l'unité \(1_E\), son image est caractérisée par
\begin{equation}\label{eq:amp-bsd-comparador-kato}
 \mathfrak c_{K/\mathrm{FERS}}(\mathfrak z_K)
 =\left[T^{-d}\rho_{E,\ell}(T)\det S_E(T)\right]_{T=0},
 \qquad
 \rho_{E,\ell}(0)=1.
\end{equation}
L'unité de \(\rho_{E,\ell}\) est orientée et vaut un à la racine ; elle ne
modifie donc ni l'ordre ni la base. Le théorème
\ref{thm:bsd-suc-bockstein} transforme \eqref{eq:amp-bsd-comparador-kato} en
\begin{equation}\label{eq:amp-bsd-kato-bockstein}
 \mathfrak c_{K/\mathrm{FERS}}(\mathfrak z_K)
 =R_{\mathrm{Boc}}^{\mathrm{der}}(S_E)
 =\tau_{\mathrm{reg}}\otimes\bigotimes_{q\geq1}\tau_q.
\end{equation}

Le comparateur de Poitou--Tate se construit page par page. Pour
\(q\geq1\), l'isomorphisme \(j_q\) de
\eqref{eq:bsd-suc-jq} et le Bockstein réduit
\(\bar\beta_q\) induisent
\begin{equation}\label{eq:amp-bsd-comparador-pt-q}
 \mathfrak c_{\mathrm{PT},q}(\tau_q)
 :=\det\!\left(j_q\circ\bar\beta_q\right)
 =\det h_E^{(q)}.
\end{equation}
La direction régulière est transportée au moyen de
\(\det\bar S_E(0)=\tau_{\mathrm{reg}}\). Le recollement de Knudsen--Mumford et les
signes de Koszul définis dans \eqref{eq:bsd-suc-rboc} fournissent l'unique
isomorphisme orienté
\begin{equation}\label{eq:amp-bsd-comparador-pt}
 \mathfrak c_{\mathrm{PT}}
 \left(\tau_{\mathrm{reg}}\otimes\bigotimes_{q\geq1}\tau_q\right)
 =\tau_{\mathrm{reg}}^{\mathrm{ht}}
  \otimes\bigotimes_{q\geq1}\det h_E^{(q)}.
\end{equation}
Sur la première page stable, la forme \(h_E^{(1)}\) est la forme de
Poincaré--Néron--Tate sur le réseau libre de Mordell--Weil ; son déterminant est
\begin{equation}\label{eq:amp-bsd-regulador-nt}
 \det h_E^{(1)}=\operatorname{Reg}(E)
\end{equation}
dans la base induite par l'unité racine.

\begin{lemma}[Déploiement de Poitou--Tate du degré arithmétique]
\label{lem:amp-bsd-grado-pt}
Soit
\[
 \operatorname{Flag}_{\mathrm{Boc}}(E)
 :=\bigoplus_{q\geq1}V_q
\]
le drapeau fini des pages positives, chacune avec la base héritée de
l'unité racine. La réduction orthogonale et les isomorphismes \(j_q\) construisent un
isomorphisme basé
\begin{equation}\label{eq:amp-bsd-psi-pt}
 \Psi_{\mathrm{PT}}:
 \bigl(E(\mathbb Q)/E(\mathbb Q)_{\mathrm{tors}}\bigr)\otimes K
 \xrightarrow{\ \sim\ }\operatorname{Flag}_{\mathrm{Boc}}(E).
\end{equation}
En particulier,
\begin{equation}\label{eq:amp-bsd-rango-bandera}
 \operatorname{rank}E(\mathbb Q)
 =\sum_{q\geq1}\dim_KV_q
 =\operatorname{ord}_T\det S_E(T).
\end{equation}
\end{lemma}

\begin{proof}
On filtre le complexe réduit par les puissances de l'idéal d'augmentation. Dans le
quotient de niveau \(q\), la direction régulière a déjà été séparée par
\(\bar S_E(0)\), et le seul bloc survivant est \(V_q\). L'équivalence
\(X^{\mathrm{orth}}\) identifie le bloc opposé à son dual, tandis que
\(j_q\bar\beta_q\) fournit l'accouplement parfait qui relève la
base du quotient au niveau précédent. En partant de la dernière page non nulle et
en répétant ce relèvement, on obtient \(\Psi_{\mathrm{PT}}\).

Dans les bases ordonnées par profondeur, la matrice de l'application résultante est
triangulaire par blocs et tous ses blocs diagonaux sont des identités : dans le
premier bloc parce que \(z_K\) et \(z_{\mathrm{PT}}\) ont la même racine, et
dans les autres parce que \(j_q\) et \(\bar\beta_q\) sont des isomorphismes basés.
Ainsi \(\Psi_{\mathrm{PT}}\) est inversible et ne perd aucune page.
En prenant les dimensions et en utilisant
\eqref{eq:bsd-suc-carry-orden}, on obtient
\eqref{eq:amp-bsd-rango-bandera}. En prenant les déterminants, le même
relèvement est exactement le recollement de
\eqref{eq:amp-bsd-comparador-pt} ; ainsi, l'égalité des rangs et l'égalité des
droites proviennent d'une seule application basée.
\end{proof}

\subsection{Triangle de localisation, rang de Smith et facteurs finis}
\label{subsec:amp-bsd-smith-sha}

Pour chaque place \(v\), la condition locale finie est un sous-complexe basé
\(\mathcal C_{E,v}^{\mathrm{fin}}\subset\mathcal C_{E,v}\), et l'on définit
\(\mathcal C_{E,v}^{\mathrm{sing}}\) par le triangle
\begin{equation}\label{eq:amp-bsd-triangulo-local}
 \mathcal C_{E,v}^{\mathrm{fin}}
 \xrightarrow{\ i_v\ }\mathcal C_{E,v}
 \longrightarrow\mathcal C_{E,v}^{\mathrm{sing}}
 \longrightarrow\mathcal C_{E,v}^{\mathrm{fin}}[1].
\end{equation}
Si \(\operatorname{res}_v\) désigne la restriction globale, l'application de
localisation est, au niveau des chaînes,
\begin{equation}\label{eq:amp-bsd-mapa-localizacion}
 \ell_E:
 \mathcal C_E^{\mathrm{glob,red}}\oplus
 \bigoplus_v\mathcal C_{E,v}^{\mathrm{fin}}
 \longrightarrow\bigoplus_v\mathcal C_{E,v},
 \qquad
 \ell_E\bigl(x,(y_v)_v\bigr)
 :=\bigl(\operatorname{res}_v x-i_vy_v\bigr)_v .
\end{equation}
La définition par cône du complexe de Selmer réduit donne le triangle
\begin{equation}\label{eq:amp-bsd-triangulo-localizacion}
 \mathcal C_E^{\mathrm{red}}
 \longrightarrow
 \mathcal C_E^{\mathrm{glob,red}}\oplus
 \bigoplus_v\mathcal C_{E,v}^{\mathrm{fin}}
 \xrightarrow{\ \ell_E\ }
 \bigoplus_v\mathcal C_{E,v}
 \longrightarrow\mathcal C_E^{\mathrm{red}}[1].
\end{equation}
L'exactitude locale--globale provient donc du cône d'une application écrite dans
\eqref{eq:amp-bsd-mapa-localizacion} ; elle n'est pas ajoutée ensuite comme une relation
entre cardinaux.

On annule dans \eqref{eq:amp-bsd-triangulo-localizacion} le facteur racine
commun, qui est le même dans les deux canaux par
\eqref{eq:bsd-suc-raiz}. Le complément résultant est noté
\begin{equation}\label{eq:amp-bsd-cono-local-reducido}
 \mathcal C_E^{\mathrm{loc,red}}
 :=q_{\mathrm{root}}\operatorname{Cone}(\ell_E)[-1].
\end{equation}
Après avoir séparé le réseau libre de Mordell--Weil et son dual au moyen de
\(X^{\mathrm{orth}}\), un choix des bases entières déjà orientées
représente ce complexe par
\begin{equation}\label{eq:amp-bsd-complejo-smith}
 \mathcal C_E^{\mathrm{loc,red}}
 \simeq[M_E\xrightarrow{\ D_E\ }N_E].
\end{equation}

\begin{lemma}[Acyclicité rationnelle et rang plein]
\label{lem:amp-bsd-aciclicidad-rango}
Le complexe de \eqref{eq:amp-bsd-complejo-smith} est acyclique après
tensorisation par \(\mathbb Q\). En particulier,
\[
 D_E\otimes\mathbb Q:M_E\otimes\mathbb Q
 \xrightarrow{\ \sim\ }N_E\otimes\mathbb Q
\]
est un isomorphisme ; \(M_E\) et \(N_E\) ont le même rang et \(D_E\) est
de rang plein.
\end{lemma}

\begin{proof}
Soit \(x\) un cycle du complexe
\(\mathcal C_E^{\mathrm{loc,red}}\otimes\mathbb Q\). La réduction a
éliminé le facteur \(p_{\mathrm{root}}\), donc
\(q_{\mathrm{root}}x=x\). Dans les bases entières orientées employées dans
\eqref{eq:amp-bsd-complejo-smith}, le complément fini--singulier se
décompose en les blocs normaux de
\cref{lem:bsd-suc-forma-normal}. Dans chaque bloc \(\lambda\), le coefficient
racine est déjà nul et le cycle s'écrit
\[
 x_{\lambda}=u_{\lambda}e_{\lambda}
                  +v_{\lambda}b_{\lambda},
 \qquad
 \partial x_{\lambda}
   =(u_{\lambda}-3c_{\lambda}v_{\lambda})g_{\lambda}=0.
\]
Comme \(g_{\lambda}\) appartient à la base libre, on a
\(u_{\lambda}=3c_{\lambda}v_{\lambda}\), et la formule explicite de la
différentielle donne
\[
 x_{\lambda}
   =v_{\lambda}(3c_{\lambda}e_{\lambda}+b_{\lambda})
   =\partial(v_{\lambda}f_{\lambda}).
\]
La somme est finie à chaque degré, de sorte que
\(x=\partial\bigl(\sum_{\lambda}v_{\lambda}f_{\lambda}\bigr)\).
Tout cycle est donc un bord par la même forme normale qui a produit
le remplisseur \eqref{eq:bsd-suc-hstar}, sans introduire
\(H_{\mathrm{FS}}\) comme prémisse. L'égalité
\(p_{\mathrm{root}}z_K=z_{\mathrm{PT}}\) est ce qui permet d'annuler la
direction racine avant ce calcul ; sans elle, il resterait une classe de degré
zéro. L'équivalence autoduale
\(X^{\mathrm{orth}}:\mathcal C_E^{\mathrm{red}}\simeq
D_3(\mathcal C_E^{\mathrm{red}})\) apparie les deux degrés restants et
transporte leurs bases avec une orientation opposée complémentaire. On en
obtient \(\operatorname{rank}M_E=\operatorname{rank}N_E\).
Enfin, l'acyclicité d'un complexe à deux termes de même rang
équivaut à l'inversibilité de \(D_E\otimes\mathbb Q\).
\end{proof}

Le lemme permet d'appliquer la forme normale de Smith sans présupposer aucun de
ses diviseurs. Il existe des changements de base orientés \(U,V\) tels que
\begin{equation}\label{eq:amp-bsd-morfismo-smith}
 U D_E V=\operatorname{diag}(s_1,\ldots,s_t),
 \qquad s_i\in\mathbb Z\setminus\{0\}.
\end{equation}
Par conséquent,
\begin{equation}\label{eq:amp-bsd-grupo-smith}
 F_E:=\operatorname{coker}D_E
 \simeq\bigoplus_{i=1}^{t}\mathbb Z/|s_i|\mathbb Z,
 \qquad
 |F_E|=\prod_{i=1}^{t}|s_i|<\infty.
\end{equation}

\begin{lemma}[Suite finie de localisation]
\label{lem:amp-bsd-exactitud-sha-tamagawa}
Le fragment de torsion de la suite longue de cohomologie de
\eqref{eq:amp-bsd-triangulo-localizacion} est
\begin{equation}\label{eq:amp-bsd-exacta-sha-tamagawa}
 0\longrightarrow\operatorname{Sha}(E)
 \xrightarrow{\ \iota_{\mathrm{Sha}}\ }F_E
 \xrightarrow{\ \partial_{\mathrm{loc}}\ }
 \bigoplus_v\Phi_v(\mathbb F_v)
 \longrightarrow0,
 \qquad c_v=|\Phi_v(\mathbb F_v)|.
\end{equation}
\end{lemma}

\begin{proof}
Un cocycle global représente une classe de
\(\operatorname{Sha}(E)\) précisément lorsque sa restriction à chaque place
admet une primitive dans
\(\mathcal C_{E,v}^{\mathrm{fin}}\). La classe de la paire formée par ce
cocycle et ses primitives locales dans le cône de
\eqref{eq:amp-bsd-mapa-localizacion} définit
\(\iota_{\mathrm{Sha}}\). Si cette classe est un bord, la composante globale est un
bord et la classe de Tate--Shafarevich est nulle ; ainsi,
\(\iota_{\mathrm{Sha}}\) est injective.

L'application \(\partial_{\mathrm{loc}}\) prend une classe du cône, la restreint à
chaque quotient
\(\mathcal C_{E,v}^{\mathrm{sing}}\) de
\eqref{eq:amp-bsd-triangulo-local} et conserve sa classe de composante.
L'identification
\[
 H^0(\mathcal C_{E,v}^{\mathrm{sing}})_{\mathrm{tors}}
 \simeq\Phi_v(\mathbb F_v)
\]
est induite par le quotient fini--singulier, non par l'ordre du groupe.
Une classe du cône a un résidu nul exactement lorsque ses représentants
locaux peuvent être choisis dans les sous-complexes finis ; d'après la définition
précédente, ce noyau est l'image de \(\iota_{\mathrm{Sha}}\).
Enfin, la dernière flèche de
\eqref{eq:amp-bsd-triangulo-local} relève chaque classe de composante au
complexe local. L'exactitude du cône
\eqref{eq:amp-bsd-triangulo-localizacion}, après retrait des deux
facteurs libres appariés par \(X^{\mathrm{orth}}\), transforme ce
relèvement en une préimage par \(\partial_{\mathrm{loc}}\). Cela démontre
la surjectivité et, par conséquent, toute
\eqref{eq:amp-bsd-exacta-sha-tamagawa}.
\end{proof}

Maintenant, et seulement après le lemme
\ref{lem:amp-bsd-aciclicidad-rango}, le groupe \(F_E\) est fini.
L'injection de \eqref{eq:amp-bsd-exacta-sha-tamagawa} démontre alors la
finitude de \(\operatorname{Sha}(E)\). En prenant les ordres dans cette suite
exacte, on obtient
\begin{equation}\label{eq:amp-bsd-cardinal-smith}
 |F_E|=|\operatorname{Sha}(E)|\prod_v c_v.
\end{equation}
C'est le comparateur
\(\mathfrak c_{\mathrm{STS}}\) : il transporte les facteurs élémentaires de Smith
vers les facteurs de Tate--Shafarevich et de Tamagawa en conservant leur ordre dans la droite,
au lieu d'insérer leurs cardinaux une fois le terme principal calculé.

La partie de torsion se construit en appliquant le déterminant au réseau de
Mordell--Weil et à son dual de Pontryagin. Les deux facteurs finis sont duaux et
apportent
\begin{equation}\label{eq:amp-bsd-comparador-torsion}
 \mathfrak c_{\mathrm{tors}}:quad
 u\longmapsto
 \frac{u}{|E(\mathbb Q)_{\mathrm{tors}}|^2}.
\end{equation}
Le comparateur archimédien s'obtient en intégrant la différentielle de Néron
orientée sur le cycle réel basé :
\begin{equation}\label{eq:amp-bsd-comparador-periodo}
 \Omega_E:=\int_{E(\mathbb R)}|\omega_E|,
 \qquad
 \mathfrak c_\infty:quad u\longmapsto\frac{u}{\Omega_E}.
\end{equation}
La différentielle et le cycle sont fixés avant d'évaluer \(L(E,s)\) ; ainsi
\eqref{eq:amp-bsd-comparador-periodo} ne normalise pas rétrospectivement la
section analytique.

\subsection{Séparation causale des primitives, des identités et des conséquences}
\label{subsec:amp-bsd-separacion-causal}

Le critère générateur par générateur du
lemme \ref{lem:amp-bsd-bases-emparejadas-kato-fers} audite à chaque degré
l'équivalence déjà construite. Il ne conditionne pas son existence. L'ordre logique
du paquet est
\begin{equation}\label{eq:amp-bsd-diagrama-causal}
\begin{array}{c}
 \begin{gathered}
  1_E,\ \Delta_{\mathrm{AW}},\ P_E^K,\ P_E^{\mathrm{PT}},
  \ h_*=-vf,\ X^{\mathrm{orth}},\\
  \{\,\mathcal C_{E,v}^{\mathrm{fin}}\to\mathcal C_{E,v}\,\}_v,\\
  \text{bases de Mordell--Weil et de Pontryagin ;}\\
  \text{différentielle et cycle de Néron}
 \end{gathered}
 \\[2mm]\Downarrow\\[-1mm]
 \begin{gathered}
  \Phi_{K/\mathrm{FERS}}\text{ équivalence},\quad
  \rho_{E,\ell}(0)=1,\quad
  p_{\mathrm{root}}z_K=z_{\mathrm{PT}},\\
  R_{\mathrm{Boc}}^{\mathrm{der}}=\operatorname{LT}_T(S_E),\quad
  \det(j_q\bar\beta_q)=\det h_E^{(q)},\\
  H^*(\mathcal C_E^{\mathrm{loc,red}}\otimes\mathbb Q)=0,\quad
  D_E\otimes\mathbb Q\text{ isomorphisme},\\
  0\to\operatorname{Sha}(E)\to F_E
  \to\bigoplus_v\Phi_v(\mathbb F_v)\to0
 \end{gathered}
 \\[2mm]\Downarrow\\[-1mm]
 \begin{gathered}
  d_{\mathrm{an}}=\operatorname{rank}E(\mathbb Q),\qquad
  \operatorname{Reg}(E),\qquad
  |\operatorname{Sha}(E)|\prod_vc_v,\\
  |E(\mathbb Q)_{\mathrm{tors}}|^{-2},\qquad
  \Omega_E^{-1},\qquad
  \text{égalité du terme principal}.
 \end{gathered}
\end{array}
\end{equation}
La première ligne contient uniquement des objets et des applications déjà basés, avec
la vérification générateur par générateur indiquée ; la deuxième, des identités
démontrées par ces applications ; la troisième, leurs conséquences sur la droite
déterminante. En particulier, aucun facteur de la dernière ligne n'est
employé pour normaliser rétroactivement une flèche de la première.

\begin{proposition}[Paquet exact de comparateurs]
\label{prop:amp-bsd-paquete-comparadores-exacto}
Les comparateurs de
\eqref{eq:amp-bsd-cadena-lineas} proviennent d'un unique paquet d'applications de
complexes et satisfont, dans l'ordre de construction,
\begin{equation}\label{eq:amp-bsd-paquete-identidades}
 \begin{gathered}
 \Phi_{K/\mathrm{FERS}}(z_K)=z_{\mathrm{Iw}},\qquad
 \operatorname{gr}^{F}\Phi_{K/\mathrm{FERS}}=I,\\
 \det\nolimits_{\mathrm{alt}}\Phi_{K/\mathrm{FERS}}
   =\rho_{E,\ell}(T),\qquad \rho_{E,\ell}(0)=1,\\
 \det(j_q\bar\beta_q)=\det h_E^{(q)}\quad(q\geq1),\\
 \Psi_{\mathrm{PT}}:
  \bigl(E(\mathbb Q)/E(\mathbb Q)_{\mathrm{tors}}\bigr)\otimes K
  \xrightarrow{\sim}\operatorname{Flag}_{\mathrm{Boc}}(E),\\
 p_{\mathrm{root}}z_K=z_{\mathrm{PT}},\qquad
 h_*=-vf,\quad \partial h_*=z_{\mathrm{PT}}-z_K,\\
 H^*(\mathcal C_E^{\mathrm{loc,red}}\otimes\mathbb Q)=0,\qquad
 U D_E V=\operatorname{diag}(s_1,\ldots,s_t),\\
 0\longrightarrow\operatorname{Sha}(E)
  \longrightarrow F_E
  \longrightarrow\displaystyle\bigoplus_v\Phi_v(\mathbb F_v)
  \longrightarrow0 .
 \end{gathered}
\end{equation}
La composition induite sur les droites déterminantes est exactement
\begin{equation}\label{eq:amp-bsd-paquete-composicion}
 \det(\Phi_{K/\mathrm{FERS}})
 \xrightarrow{\ \det(j_q\bar\beta_q)\ }
 \det(D_E)
 \xrightarrow{\ \mathrm{Smith}\ }
 \frac{|\operatorname{Sha}(E)|\prod_vc_v}
      {|E(\mathbb Q)_{\mathrm{tors}}|^2\,\Omega_E},
\end{equation}
avec les signes de Knudsen--Mumford et les orientations de
\eqref{eq:amp-bsd-cadena-lineas}. En particulier, le membre droit de
\eqref{eq:amp-bsd-paquete-composicion} est la coordonnée finale de l'élément
transporté et non une donnée servant à construire l'une des flèches.
\end{proposition}

\begin{proof}
La première ligne de \eqref{eq:amp-bsd-paquete-identidades} est
\eqref{eq:amp-bsd-phi-kato-fers} et
\eqref{eq:amp-bsd-matriz-kato-fers}--
\eqref{eq:amp-bsd-rho-determinantal}. La deuxième est
\eqref{eq:amp-bsd-comparador-pt-q} et le lemme
\ref{lem:amp-bsd-grado-pt}. La troisième combine l'égalité de la racine avec
le remplisseur explicite \eqref{eq:bsd-suc-hstar} ; son extension linéaire dans la
forme normale induit le lecteur auxiliaire utilisé dans le lemme
\ref{lem:amp-bsd-aciclicidad-rango}. En se restreignant au complément de la
racine, cette réduction contracte tout cycle rationnel. La dernière ligne est
\eqref{eq:amp-bsd-morfismo-smith} et
\eqref{eq:amp-bsd-exacta-sha-tamagawa}. Appliquer le foncteur déterminant à
ces identités, dans ce même ordre, donne
\eqref{eq:amp-bsd-paquete-composicion}. Aucune égalité ne s'obtient
en comparant d'abord les deux scalaires finaux.
\end{proof}

\begin{lemma}[Indépendance causale du terme principal]
\label{lem:amp-bsd-independencia-termino-principal}
Soit \(\mathscr P_E\) l'ensemble de données
\begin{equation}\label{eq:amp-bsd-primitivas-paquete}
 \mathscr P_E:=
 \left(
  1_E,\Delta_{\mathrm{AW}},P_E^K,P_E^{\mathrm{PT}},
  h_*=-vf,X^{\mathrm{orth}},
  \{\mathcal C_{E,v}^{\mathrm{fin}}\xrightarrow{i_v}
     \mathcal C_{E,v}\}_v,
  \text{bases et orientations},
  \omega_E,\gamma_E^{\mathbb R}
 \right),
\end{equation}
où \(\omega_E\) est la différentielle de Néron et
\(\gamma_E^{\mathbb R}\) le cycle réel orienté. Toutes les applications du
paquet de la proposition
\ref{prop:amp-bsd-paquete-comparadores-exacto} sont déterminées par
\(\mathscr P_E\). Ne font pas partie de \(\mathscr P_E\) les numéraux
\[
 L^{(r)}(E,1),\quad r,\quad\operatorname{Reg}(E),\quad
 |\operatorname{Sha}(E)|,\quad(c_v)_v,\quad
 |E(\mathbb Q)_{\mathrm{tors}}|,\quad\Omega_E.
\]
Ces valeurs s'obtiennent, respectivement, en évaluant la section analytique,
en prenant le degré du drapeau, en calculant des déterminants et des formes de Smith,
en prenant les ordres de groupes finis et en intégrant
\(\omega_E\) sur \(\gamma_E^{\mathbb R}\).
\end{lemma}

\begin{proof}
La formule \eqref{eq:amp-bsd-phi-kato-fers} utilise uniquement
\(1_E,\Delta_{\mathrm{AW}},P_E^K\) et la retenue. Les applications de Bockstein et de
Poitou--Tate utilisent la filtration, \(P_E^{\mathrm{PT}}\), les bases et
l'autodualité. Le cône de \eqref{eq:amp-bsd-mapa-localizacion} utilise les
applications \(i_v\) et les restrictions locales ; sa contraction et sa forme de
Smith utilisent le remplisseur explicite \eqref{eq:bsd-suc-hstar}, son extension
linéaire auxiliaire, \(X^{\mathrm{orth}}\) et les bases entières. Les deux
derniers comparateurs utilisent les réseaux de torsion et la
paire \((\omega_E,\gamma_E^{\mathbb R})\). Il s'agit d'une énumération exhaustive
des arguments qui figurent dans les formules de construction. Les
numéraux énumérés dans l'énoncé n'apparaissent qu'après application de la
cohomologie, du déterminant, du cardinal ou de l'intégration, de sorte qu'ils ne peuvent pas
sélectionner rétrospectivement les applications.
\end{proof}

\subsection{Composition, degrés et terme principal}
\label{subsec:amp-bsd-composicion}

Définissons le comparateur total par la composition dans l'ordre écrit :
\begin{equation}\label{eq:amp-bsd-comparador-total}
 \mathfrak C_E:=
 \mathfrak c_\infty\circ
 \mathfrak c_{\mathrm{tors}}\circ
 \mathfrak c_{\mathrm{STS}}\circ
 \mathfrak c_{\mathrm{PT}}\circ
 \mathfrak c_{K/\mathrm{FERS}}.
\end{equation}
Comme \(p_{\mathrm{root}}z_K=z_{\mathrm{PT}}\), toutes les flèches reçoivent la
même unité basée. Le remplisseur \(h_*=-vf\) de
\eqref{eq:bsd-suc-hstar} remplit explicitement la différence ; son extension
linéaire dans la base normale produit le lecteur homotopique auxiliaire et démontre
qu'aucune classe supplémentaire ne modifie l'élément de la droite pendant le
transport. Le lecteur n'est pas une prémisse distincte.

\begin{theorem}[Construction des comparateurs basés]
\label{thm:amp-bsd-comparadores-construidos}
Pour toute courbe elliptique \(E/\mathbb Q\) munie des complexes, des bases,
des orientations et des données locales définis précédemment, les applications
\eqref{eq:amp-bsd-phi-kato-fers},
\eqref{eq:amp-bsd-comparador-pt},
\eqref{eq:amp-bsd-triangulo-localizacion},
\eqref{eq:amp-bsd-comparador-torsion} et
\eqref{eq:amp-bsd-comparador-periodo}, avec les lemmes
\ref{lem:amp-bsd-aciclicidad-rango} et
\ref{lem:amp-bsd-exactitud-sha-tamagawa}, construisent les cinq comparateurs de
\eqref{eq:amp-bsd-cadena-lineas}. Leur composition préserve l'unité, le degré
et l'orientation, et satisfait
\begin{equation}\label{eq:amp-bsd-igualdad-grados}
 d_{\mathrm{an}}=\operatorname{ord}_T\det S_E(T)
 =d_{\mathrm{ar}}=\operatorname{rank}E(\mathbb Q).
\end{equation}
De plus, \(\operatorname{Sha}(E)\) est fini et l'égalité de l'élément
distingué dans \(\mathcal L_E\) se publie sous la forme
\begin{equation}\label{eq:amp-bsd-termino-principal}
 \frac{L^{(r)}(E,1)}{r!\,\Omega_E}
 =\frac{|\operatorname{Sha}(E)|\,\operatorname{Reg}(E)
        \prod_v c_v}
       {|E(\mathbb Q)_{\mathrm{tors}}|^2},
 \qquad r=\operatorname{rank}E(\mathbb Q).
\end{equation}
\end{theorem}

\begin{proof}
Le comparateur Kato--FERS et la condition
\(\rho_{E,\ell}(0)=1\) identifient l'ordre analytique à
\(d=\operatorname{ord}_T\det S_E(T)\) et le terme initial à
\(R_{\mathrm{Boc}}^{\mathrm{der}}(S_E)\). L'identité
\eqref{eq:bsd-suc-lt} conserve la page régulière et chaque page positive. Les
isomorphismes \(j_q\) transportent ces pages vers les hauteurs dérivées. Le
lemme \ref{lem:amp-bsd-grado-pt} construit l'isomorphisme du drapeau avec le
réseau libre de Mordell--Weil et, avec
\eqref{eq:bsd-suc-carry-orden}, démontre
\eqref{eq:amp-bsd-igualdad-grados}.

Le lemme \ref{lem:amp-bsd-aciclicidad-rango} déduit l'acyclicité rationnelle
directement du remplisseur explicite \eqref{eq:bsd-suc-hstar}, de sa réduction
linéaire du complément, de l'annulation de la racine commune et de
\(X^{\mathrm{orth}}\) ; on en obtient le rang plein de \(D_E\), non
comme hypothèse supplémentaire. La forme de Smith
\eqref{eq:amp-bsd-morfismo-smith} produit alors le groupe fini \(F_E\).
Le lemme \ref{lem:amp-bsd-exactitud-sha-tamagawa} extrait des triangles
\eqref{eq:amp-bsd-triangulo-local} et
\eqref{eq:amp-bsd-triangulo-localizacion} la suite exacte finie ; ce n'est
qu'ensuite que l'on déduit la finitude de \(\operatorname{Sha}(E)\) et
\eqref{eq:amp-bsd-cardinal-smith}. Enfin,
\eqref{eq:amp-bsd-regulador-nt},
\eqref{eq:amp-bsd-comparador-torsion} et
\eqref{eq:amp-bsd-comparador-periodo} publient, respectivement, le régulateur,
le quotient de torsion et la période. L'égalité racine élimine toute unité
relative entre les deux publications. En évaluant le même élément basé aux
deux extrémités de \eqref{eq:amp-bsd-comparador-total}, on obtient
\eqref{eq:amp-bsd-termino-principal}.
\end{proof}

\subsection{Prémisses et critères de réfutation}
\label{subsec:amp-bsd-falsadores}

La construction part de l'unité généalogique \(1_E\) et de la diagonale
d'Alexander--Whitney. Elle utilise également les
publications couplées, tous les
Bockstein et la page régulière, le remplisseur explicite \(h_*=-vf\),
l'autodualité réduite \(X^{\mathrm{orth}}\), les applications de chaînes du
triangle de localisation et les bases orientées de torsion et de période. L'équivalence
Kato--FERS et sa normalisation à un, le rang plein de \(D_E\),
l'exactitude de \eqref{eq:amp-bsd-exacta-sha-tamagawa} et la finitude de
\(\operatorname{Sha}(E)\) sont des conclusions des lemmes précédents, non des
prémisses indépendantes. Le critère générateur par générateur du lemme
\ref{lem:amp-bsd-bases-emparejadas-kato-fers} est un contrôle non directeur
qui audite localement le premier comparateur. Les critères qui réfuteraient une flèche précise
sont :
\begin{enumerate}
 \item un coefficient racine différent de un, qui introduirait une unité
       relative entre \(z_K\) et \(z_{\mathrm{PT}}\) ;
 \item un \(j_q\) non inversible ou l'omission d'une page de Bockstein, qui
       modifierait le degré ou le déterminant de hauteur ;
 \item un diviseur de Smith nul, qui empêcherait la finitude de \(F_E\) ;
 \item l'échec de l'exactitude de
       \eqref{eq:amp-bsd-exacta-sha-tamagawa}, qui romprait l'identification
       des facteurs finis ;
 \item une inversion d'orientation dans la torsion ou la période, qui changerait
       l'élément basé même si elle préservait sa droite non orientée.
\end{enumerate}
La projection de Manin sans histoire et le changement d'échelle visible modulo trois
violent le premier critère ; ils constituent donc des preuves de nécessité pour
ces modèles réduits et non des prémisses du théorème
\ref{thm:amp-bsd-comparadores-construidos}.
